% Source: https://artofproblemsolving.com/community/c3h2495359_ratios__a_shape
% Stable identity: aops-topic-2495359-q01
\textbf{(a)} The circumferences are $2\pi\cdot3=6\pi$ and $2\pi\cdot5=10\pi$, in the ratio $\boxed{3:5}$. The areas are $9\pi$ and $25\pi$, in the ratio $\boxed{9:25}$.

\textbf{(b)} The perimeters are $2(2+3)=10$ and $2(6+9)=30$, in the ratio $\boxed{1:3}$. The areas are $6$ and $54$, in the ratio $\boxed{1:9}$.

\begin{asy}
usepackage("amsmath");
size(6cm);
defaultpen(fontsize(10pt));
draw(circle((0,0),3),linewidth(0.8)); dot((0,0)); draw((0,0)--(3,0)); label("$3$",(1.5,0),N);
draw(circle((10,0),5),linewidth(0.8)); dot((10,0)); draw((10,0)--(15,0)); label("$5$",(12.5,0),N);
draw(box((-1.5,-12),(1.5,-10)),linewidth(0.8)); label("$3$",(0,-12),S); label("$2$",(1.5,-11),E);
draw(box((5.5,-18),(14.5,-12)),linewidth(0.8)); label("$9$",(10,-18),S); label("$6$",(14.5,-15),E);
\end{asy}

\textbf{(c)} A circle of radius $r$ becomes a circle of radius $kr$: its circumference $2\pi(kr)=k\cdot2\pi r$ is multiplied by $k$, and its area $\pi(kr)^2=k^2\cdot\pi r^2$ by $k^2$. A rectangle $a\times b$ becomes $ka\times kb$: its perimeter $2(ka+kb)=k\cdot2(a+b)$ is multiplied by $k$, and its area $(ka)(kb)=k^2\cdot ab$ by $k^2$.

So for two shapes of the same kind whose lengths are in the ratio $1:k$, perimeters are in the ratio $1:k$ and areas in the ratio $1:k^2$. In (a), $k=\dfrac53$ and $9:25=3^2:5^2$; in (b), $k=3$ and $1:9=1^2:3^2$.
